\shorthandoff{:!} %règle un problème de compatibilité avec babel
\begin{tikzpicture}[scale=\scaleFigure]
	%ecran oscilloscope
	\draw[black,thick,rounded corners=10pt] (0,0) rectangle (10,8) (0.5,-0.5) node[right]{\textcolor{Black}{\scriptsize{base de temps : 
	\SI{20}{\micro\second/division}}}}
	(0.5,-1.1) node[right]{\textcolor{Black}{\scriptsize{calibre vertical : \SI{1}{V/division}}}} ;
	\begin{scope} % permet d'appliquer des propriétés à une partie de code seulement
		\clip (-2,0.05) rectangle (9.95,7.95);%créer une découpe de la fenêtre comme indiqué
		\foreach \x in {1,2,...,9}%Démarrage d'une boucle : "pour tout x appartenant à ..."
		{
			%la boucle est délimitée par une paire d'accolade
			\draw[gray!20, very thin] (\x,0.05) -- (\x,7.95);%thin = fin, il existe aussi very thin et ultra thin
		}
		\foreach \y in {1,2,...,7}
		{
			\draw[gray!20, very thin] (0.05,\y) -- (9.95,\y);
		}
		\foreach \x in {0.2,0.4,...,9.8}
		{
			\draw[gray!20, very thin] (\x,4.1) -- (\x,3.9);
		}
		\foreach \y in {0.2,0.4,...,7.8}
		{
			\draw [gray!20, very thin] (5.1,\y) -- (4.9,\y);
		}
		%FIN ecran oscilloscope
		% Les tensions d'entrée et de sortie
		\draw[yshift=4cm,color=black,thick,domain=0:10, smooth, samples=150] plot ({\x},{-0.5+2.5*sin(1.256*\x r+pi/10 r)});
   	     \draw (-1.7,1) node [right]  {\scriptsize{0V}};
		%legende
		%un triangle définition de la ligne de masse
		\fill [black] (-0.5,0.75) -- ++(30:0.5) -- ++(150:0.5) --cycle ;
	
	\end{scope}
\end{tikzpicture} 
\shorthandon{:!}